\clearpage
\begin{proofbox}
\textbf{统一前提与选法说明。}

记 $k=b/a>0$，$c^2=a^2+b^2$，$e=c/a>1$。
由关于 $x$ 轴及 $y$ 轴的对称性，\textbf{只证}右焦点
$F=(c,0)$、过 $F$ 作斜率 $k$ 的直线
\[
\ell:\ y=k(x-c).
\]
设 $N$ 为此线与另一条渐近线 $y=-kx$ 的交点，
$M$ 为此线与右支双曲线的交点。
下面提供\textbf{最快、最低门槛、常规}三种证明，
并给出斜坐标补充解释；前三法均可独立完成。

\medskip
\textbf{方法一：最快证明【双曲线因式分解＋分点线性量】。}

\begin{center}\begin{tikzpicture}[font=\small,>=stealth]
\draw[gray!45,rounded corners=3pt] (0,0) rectangle (14.2,9.5);
\node at (7.1,8.85) {方法一：三条平行等值线，把截线分成对应的两段};

\begin{scope}[shift={(.9,5.9)},scale=.85]
\draw[violet,dashed] (0,0)--(3.3,-4.4);
\draw[violet,dashed] (0,2.4)--(4.8,-4);
\draw[violet,dashed] (3,8/3)--(6.2,-1.6);

\coordinate (F) at (5,0);\coordinate (M) at (3.4,-32/15);\coordinate (N) at (2.5,-10/3);
\draw[blue!75!black,very thick] (F)--(M);
\draw[orange!85!black,very thick] (M)--(N);
\foreach \q in {F,M,N}\fill (\q) circle(2.8pt);
\node[above right=4pt] at (F) {$F$};\node[right=5pt] at (M) {$M$};\node[below left=4pt] at (N) {$N$};

\node[anchor=west,violet] at (3.35,-4.4) {$L=0$};
\node[anchor=west,violet] at (4.85,-4) {$L=a^2/c$};
\node[anchor=west,violet] at (6.25,-1.6) {$L=c$};
\end{scope}
\node[anchor=west,align=left] at (8.6,6.8) {$L=x+\dfrac ab y$\\[14pt]平行等值线\\[7pt]对应同一直线的分段};
\node[anchor=west,align=left] at (8.6,3.5) {蓝段的 $L$ 差：\\[8pt]$c-a^2/c$\\[10pt]橙段的 $L$ 差：\\[8pt]$a^2/c-0$};
\node at (7.1,.65) {图取 $a=3,b=4,c=5$；$L(F),L(M),L(N)$ 依次为 $5,9/5,0$。};

\end{tikzpicture}
\end{center}

\textbf{第1步：把曲线写成两个渐近线因子的乘积。}
\[
\frac{x^2}{a^2}-\frac{y^2}{b^2}=1
\quad\Longleftrightarrow\quad
\left(x-\frac ab y\right)\left(x+\frac ab y\right)=a^2.
\]
$\ell$ 上恒有 $x-\tfrac ab y=c$；因此\textbf{在交点 $M$ 处}
\[
\boxed{x_M+\frac ab y_M=\frac{a^2}{c}}.
\]

\textbf{第2步：找一个沿直线均匀变化的量。}
令 $L(x,y)=x+\tfrac ab y$。图中 $L$ 的等值线互相平行；沿截线的分段比，等于对应 $L$ 的变化量之比（平行线分线段成比例）。
\[
L(F)=c,\qquad L(M)=\frac{a^2}{c},\qquad
L(N)=0
\quad\bigl(N\in y=-kx\bigr).
\]
故线段比\textbf{直接等于} $L$ 的变化量之比：
\[
\begin{aligned}
\frac{FM}{MN}
&=\frac{L(F)-L(M)}{L(M)-L(N)}\\[2pt]
&=\frac{c-a^2/c}{a^2/c}
=\frac{c^2-a^2}{a^2}
=\boxed{\frac{b^2}{a^2}=e^2-1}.
\end{aligned}
\]
由 $\overrightarrow{FM}=\lambda\overrightarrow{MN}$（$\lambda>0$），
两向量同向且 $\lambda=FM/MN$，所以
\[
\boxed{e=\sqrt{1+\lambda}}.
\]
\textbf{为什么最快：}无需算 $M,N$ 的坐标，无需解二次方程，
用因式分解把曲线交点变成“一次量”。
\textbf{适用前提：}须理解线性函数沿直线的分点比；新手可先用方法二。
\end{proofbox}

\clearpage
\begin{proofbox}
\textbf{方法二：最低门槛【只设 $t$＋平方差；首推初学者】。}

\begin{center}\begin{tikzpicture}[font=\small,>=stealth]
\draw[gray!45,rounded corners=3pt] (0,0) rectangle (14.2,9.5);
\node at (7.1,8.85) {方法二：从 F 向左走 t，再向下走 kt，便到达 M};

\begin{scope}[shift={(.3,5.8)},scale=1.05]

\draw[teal!75!black,thick,samples=90,domain=-3:1.6,variable=\v] plot ({3*sqrt(1+\v*\v/16)},{\v});

\coordinate (F) at (5,0);\coordinate (M) at (3.4,-32/15);\coordinate (N) at (2.5,-10/3);
\draw[blue!75!black,very thick] (F)--(M);
\draw[orange!85!black,very thick] (M)--(N);
\foreach \q in {F,M,N}\fill (\q) circle(2.8pt);
\node[above right=4pt] at (F) {$F$};\node[right=5pt] at (M) {$M$};\node[above left=5pt] at (N) {$N$};

\draw[blue!75!black,dashed,->] (F)--(3.4,0) node[midway,above=4pt] {$t$};
\draw[blue!75!black,dashed,->] (3.4,0)--(M) node[midway,right=4pt] {$kt$};
\draw[gray] (3.4,.2)--(3.4,0)--(5.6,0);
\draw[gray,dotted] (N)--(2.5,-4.2) (M)--(3.4,-4.2) (F)--(5,-4.2);
\draw[<->,blue!75!black] (3.4,-3.9)--(5,-3.9) node[midway,below] {$t$};
\draw[<->,orange!85!black] (2.5,-3.9)--(3.4,-3.9);
\node[orange!85!black,anchor=east] at (2.35,-3.9) {$c/2-t$};
\end{scope}
\node[anchor=west,align=left] at (8.4,6.3) {$F=(c,0)$\\[12pt]$M=(c-t,-kt)$\\[12pt]$N_x=c/2$};
\node[anchor=west,align=left] at (8.4,3.2) {$\displaystyle\frac{FM}{MN}=\frac{t}{c/2-t}$\\[12pt]图中 $t=8/5<c/2$};
\node at (7.1,.65) {先由方向确定设点；再用双曲线方程求 $t$。};

\end{tikzpicture}
\end{center}

\textbf{第1步：先求容易求的点 $N$。}
\[
\begin{cases}
y=k(x-c),\\
y=-kx
\end{cases}
\quad\Longrightarrow\quad
\boxed{N=\left(\frac c2,-\frac{kc}{2}\right)}.
\]
图中两条直线先确定 $N_x=c/2$。

\textbf{第2步：沿已知方向，只设一个数 $t$。}
依照图中的水平、竖直移动，设向左位移为 $t$（求解后核查 $t>0$）：
\[
\boxed{M=(c-t,-kt)}.
\]
把点代入双曲线并利用 $k=b/a$：
\[
\frac{(c-t)^2}{a^2}-\frac{k^2t^2}{b^2}=1
\quad\Longrightarrow\quad
\frac{(c-t)^2-t^2}{a^2}=1.
\]
展开\textbf{平方差}即可，二次项完全消失：
\[
c^2-2ct=a^2
\quad\Longrightarrow\quad
\boxed{t=\frac{c^2-a^2}{2c}=\frac{b^2}{2c}}.
\]
因为 $0<b^2<c^2$，故 $0<t<c/2$，
从而 $F-M-N$ 的位置次序得到严格确认，而非仅凭图猜测。

\textbf{第3步：同一直线，横向差之比就是长度比。}
\[
\frac{FM}{MN}
=\frac{c-(c-t)}{(c-t)-c/2}
=\frac{t}{c/2-t}
=\frac{b^2/(2c)}{(c^2-b^2)/(2c)}
=\boxed{\frac{b^2}{a^2}}.
\]
再由 $c^2=a^2+b^2$，得到
\[
\lambda=\frac{b^2}{a^2},\qquad
\boxed{e=\sqrt{1+\lambda}}.
\]

\textbf{门槛低在哪里：}无需斜坐标、向量技巧或参数方程；
只用\textbf{两条直线联立、一个 $t$、平方差、共线线段比}。
\end{proofbox}

\clearpage
\begin{proofbox}
\textbf{方法三：常规解析几何【写直线＋联立求交点＋距离公式】。}

\begin{center}\begin{tikzpicture}[font=\small,>=stealth]
\draw[gray!45,rounded corners=3pt] (0,0) rectangle (14.2,9.5);
\node at (7.1,8.85) {方法三：坐标作垂足，横向差乘同一系数就是实际长度};

\begin{scope}[shift={(.3,6.1)},scale=1.05]

\draw[teal!75!black,thick,samples=90,domain=-4.2:1.6,variable=\v] plot ({3*sqrt(1+\v*\v/16)},{\v});

\coordinate (F) at (5,0);\coordinate (M) at (3.4,-32/15);\coordinate (N) at (2.5,-10/3);
\draw[blue!75!black,very thick] (F)--(M);
\draw[orange!85!black,very thick] (M)--(N);
\foreach \q in {F,M,N}\fill (\q) circle(2.8pt);
\node[above right=4pt] at (F) {$F$};\node[right=5pt] at (M) {$M$};\node[below left=4pt] at (N) {$N$};

\draw[gray,->] (1.4,0)--(6,0) node[right] {$x$};
\draw[gray,dashed] (M)--(3.4,0) (N)--(2.5,0);
\draw[blue!75!black,thick] (3.4,.5)--(5,.5);
\draw[orange!85!black,thick] (2.5,.5)--(3.4,.5);
\node[above=5pt] at (5,.5) {$c$};
\node[above=5pt] at (3.4,.5) {$x_M$};
\node[above left=5pt] at (2.5,.5) {$c/2$};
\draw[blue!75!black,->] (5.2,-1.8)--(4.35,-.87);
\node[blue!75!black,anchor=west] at (5.25,-1.85) {$FM$};
\end{scope}
\node[anchor=west,align=left] at (8.3,6.5) {横向差：\\[10pt]蓝 $b^2/(2c)$\\[10pt]橙 $a^2/(2c)$};
\node[anchor=west,align=left] at (8.3,3.5) {实际长度都乘\\[12pt]$\sqrt{1+k^2}=c/a$};
\node at (7.1,.65) {图取 $a=3,b=4,c=5$；横向差 $8/5,9/10$，实长 $8/3,3/2$。};

\end{tikzpicture}
\end{center}

\textbf{第1步：求 $N$。}
把 $y=k(x-c)$ 与 $y=-kx$ 联立，得
\[
\boxed{N=\left(\frac c2,-\frac{bc}{2a}\right)}.
\]

\textbf{第2步：求 $M$。}
将 $y=\tfrac ba(x-c)$ 代入 $H$：
\[
\frac{x^2}{a^2}-\frac{(x-c)^2}{a^2}=1
\quad\Longrightarrow\quad
2cx-c^2=a^2.
\]
所以
\[
\boxed{x_M=\frac{a^2+c^2}{2c}
=c-\frac{b^2}{2c}},
\]
\[
\boxed{y_M=\frac ba(x_M-c)=-\frac{b^3}{2ac}}.
\]
并且 $a<x_M<c$：
\[
x_M-a=\frac{(c-a)^2}{2c}>0,
\qquad c-x_M=\frac{b^2}{2c}>0.
\]
这是右支上的有效交点；代回直线满足题设。

\textbf{第3步：求线段长。}
图中两段的水平投影分别为 $c-x_M$、$x_M-c/2$；由勾股定理，两段实际长度都乘 $\sqrt{1+k^2}=c/a$。
因此
\[
\begin{aligned}
FM&=(c-x_M)\frac ca
=\frac{b^2}{2c}\frac ca=\boxed{\frac{b^2}{2a}},\\[3pt]
MN&=(x_M-c/2)\frac ca
=\frac{a^2}{2c}\frac ca=\boxed{\frac a2},\\[3pt]
FN&=(c-c/2)\frac ca=\boxed{\frac{c^2}{2a}}.
\end{aligned}
\]
故
\[
\boxed{\frac{FM}{MN}=\frac{b^2}{a^2}
=e^2-1=\lambda},\quad
\boxed{e=\sqrt{1+\lambda}}.
\]
\textbf{为什么保留此法：}是考试中最标准、最容易检查步骤的解法；
还额外得到三段精确长度。
\end{proofbox}

\clearpage
\begin{proofbox}
\textbf{方法四：拓展理解【把双曲线化为 $UV=1$ 的斜坐标法】。}

\textbf{补充原因：}如果想弄清“为什么平行渐近线使难题变简单”，
这个坐标变化能把几何原因直观地显露出来；
但它与方法一同源，\textbf{不应当算作完全不同的定理}。

\begin{center}\begin{tikzpicture}[font=\small,>=stealth]
\draw[gray!45,rounded corners=3pt] (0,0) rectangle (14.2,9.5);
\node at (7.1,8.85) {方法四：在新坐标平面中，固定 U，沿 V 轴方向读分段};

\node[anchor=west,align=left] at (.8,6.8) {$U=x/a-y/b$\\[12pt]$V=x/a+y/b$\\[16pt]双曲线变为 $UV=1$};
\node[anchor=west,align=left] at (.8,3.5) {$U=e$ 时，三点的 $V$：\\[12pt]$F:e,\quad M:1/e,\quad N:0$};
\begin{scope}[shift={(8,2)},scale=2]
\draw[gray,->] (0,0)--(2.65,0) node[right] {$U$};
\draw[gray,->] (0,0)--(0,2.6) node[above] {$V$};
\draw[teal!75!black,thick,samples=90,domain=.42:2.45] plot (\x,{1/\x});
\draw[gray,dashed] (5/3,0)--(5/3,2.4);
\draw[blue!75!black,very thick] (5/3,5/3)--(5/3,.6);
\draw[orange!85!black,very thick] (5/3,.6)--(5/3,0);
\draw[gray,dotted] (0,5/3)--(5/3,5/3) (0,.6)--(5/3,.6);
\fill (5/3,5/3) circle(1.3pt) node[right=5pt] {$F$};
\fill (5/3,.6) circle(1.3pt) node[above right=6pt] {$M$};
\fill (5/3,0) circle(1.3pt) node[below=5pt] {$N$};
\node[left=4pt] at (0,5/3) {$e$};\node[left=4pt] at (0,.6) {$1/e$};
\node[above right] at (5/3,2.4) {$U=e$};
\end{scope}
\node at (7.1,.65) {图取 $e=5/3$；这是新坐标图，原图长度不照搬，固定 $U$ 后的分段比保留。};

\end{tikzpicture}
\end{center}

定义无量纲量
\[
U=\frac xa-\frac yb,\qquad V=\frac xa+\frac yb.
\]
双曲线立即化为
\[
\boxed{UV=1}.
\]
过右焦点、斜率 $b/a$ 的直线满足
\[
\boxed{U=\frac ca=e}.
\]
于是同一直线上三个点的 $V$ 坐标分别是
\[
V_F=e,\qquad V_M=\frac1e,\qquad V_N=0.
\]
图中竖线 $U=e$ 上直接读出两个 $V$ 差。变换回原图时，因 $U$ 固定，长度只乘一个公共系数，所以分段比保持：
\[
\frac{FM}{MN}
=\frac{V_F-V_M}{V_M-V_N}
=\frac{e-1/e}{1/e}
=\boxed{e^2-1}.
\]
\textbf{优点：}几乎只有三个数 $e,1/e,0$；非常适合记忆本质。
\textbf{代价：}要先理解斜坐标变换，对只学过普通直角坐标的学生不是首选。

\medskip
\textbf{全构型对称性与边界核查。}
上文取右焦点及正斜率，另一种斜率经关于 $x$ 轴反射得到；
换成左焦点再关于 $y$ 轴反射得到。
反射保持线段长度及同向向量倍数，故结论对四种选择均成立。
并且 $a,b>0$ 保证 $e>1$、$\lambda>0$。

a) 若题目把向量倒过来写为
$\overrightarrow{MN}=\mu\overrightarrow{FM}$（$\mu>0$），
则 $\mu=1/\lambda$，必须改用 $e=\sqrt{1+1/\mu}$。

b) 若只是“过焦点任意直线”，不保证平行渐近线，
则关键因式不再恒定，不能调用本结论。

\medskip
\textbf{选法建议：}
\emph{考场最快}用方法一；\emph{最低认知负担}用方法二；
\emph{规范答卷和求绝对长度}用方法三；
\emph{理解公式为什么整齐}用方法四。
\end{proofbox}
